\section{D-MIMO Power Model}
\label{sec:dmimo_pwr}

The model of \Cref{sec:pwr_model} describes one central site: a single co-located array
whose antennas, converters, and processors share a physical location, a clock, and a power
supply. The distributed system of \Cref{sec:dmimo_sysmodel} replaces that site by
$L$ \glspl{ap} of $M$ antennas each, spread over the coverage area and tied to a
\gls{cpu} by fronthaul links. Three things change as a result. First, the hardware
blocks become per-AP quantities that must be summed over the network, each at its
own operating point. Second, the transport between the \glspl{ap} and the \gls{cpu} has its own power consumption, wich is not modeled in the single-site model. Third, the digital
processing chain is split across two locations (\gls{ap} and \gls{cpu}). This requires design choices about where a given operation runs. This section extends the parametric model along those three axes while keeping the frame-averaging machinery of \Cref{sec:frame} intact, so that the single-site model is recovered as a special case by setting $L=1$ with no fronthaul and no \gls{cpu}.

Throughout, $l=1,\dots,L$ indexes APs, $k=1,\dots,K$ users, $q=1,\dots,Q$ subcarriers, and $i\in\{\mathrm{DL},\mathrm{UL}\}$ the link direction, matching \Cref{sec:dmimo_sysmodel}. Hardware parameters carried over from \Cref{sec:pwr_model} keep their symbols and now refer to a single AP.

\subsection{Network power decomposition}
\label{ssec:dmimo_structure}

Let $\mathcal{S}\subseteq\{1,\dots,L\}$ be the set of APs that are switched on;
an AP outside $\mathcal{S}$ is in deep sleep and draws a fraction
$\delta^{\mathrm{ds}}\in[0,1]$ of its base consumption
$P^{0}_{\mathrm{AP},l}$. The total distributed network consumption is
\begin{align}
\label{eq:pnet}
P_{\mathrm{net}} ={}& \sum_{l\in\mathcal{S}}
\left(P_{\mathrm{AP},l} + \frac{\overline{P_{\mathrm{FH},l}}}{\eta_{\mathrm{FH,sc}}}\right)
\nonumber\\
&+ \sum_{l\notin\mathcal{S}} \delta^{\mathrm{ds}}
\left(P^{0}_{\mathrm{AP},l} + \frac{P^{0}_{\mathrm{FH},l}}{\eta_{\mathrm{FH,sc}}}\right)
+ \frac{\overline{P_{\mathrm{CPU}}}}{\eta_{\mathrm{CPU,sc}}},
\end{align}
where the consumption of an active AP retains the three-block structure of
\eqref{eq:pcons},
\begin{equation}
\label{eq:pap}
P_{\mathrm{AP},l} =
\frac{\overline{P_{\mathrm{dig},l}}}{\eta_{\mathrm{dig,sc}}}
+ \frac{\overline{P_{\mathrm{ana},l}}}{\eta_{\mathrm{ana,sc}}}
+ \frac{\overline{P_{\mathrm{PA},l}}}{\eta_{\mathrm{PA,sc}}},
\end{equation}
and $\eta_{\mathrm{FH,sc}},\eta_{\mathrm{CPU,sc}}\in(0,1]$ are the
supply-and-cooling efficiencies of the fronthaul transceivers and of the central
unit. The CPU efficiency plays the role of an inverse power-usage effectiveness:
the central unit sits in a room with its own cooling, which is not the same
environment as an outdoor AP enclosure, so $\eta_{\mathrm{CPU,sc}}$ should not be
set equal to the $\eta_{c,\mathrm{sc}}$ values fitted for a macro base station. \todo{find proper values, look at power usage effectiveness (PUE) of datacenters}
Setting $L=1$, $\mathcal{S}=\{1\}$, $M=M_{\mathrm{ant}}$, and dropping the
fronthaul and CPU terms returns \eqref{eq:pcons} unchanged.

The set $\mathcal{S}$ determines how many \glspl{ap} serve the network. As a first step we assume that all $L$ \glspl{ap} are active. 

\todo{see if we want to keep the next paragraph:}
The set $\mathcal{S}$ is the first genuinely new degree of freedom. A co-located
array can shut down individual chains but not the site; a distributed deployment
can shut down whole APs while the network keeps serving, and the coverage
redundancy of cell-free operation is what makes that possible. AP selection
therefore acts directly on \eqref{eq:pnet}, and \Cref{rem:fh_ceiling} shows that
it also relieves the fronthaul.

\subsection{Serving clusters, per-AP load, and the shared frame}
\label{ssec:dmimo_frame}

In the fully distributed operation of \Cref{sec:dmimo_sysmodel} every AP serves
every user. The power model should not assume this, because restricting each AP
to a subset of users is the standard way of keeping cell-free operation scalable
and is also the main lever on fronthaul traffic. Let
$\mathcal{D}_{l}\subseteq\{1,\dots,K\}$ be the set of users served by AP $l$ and
$\mathcal{L}_{k}=\{l:k\in\mathcal{D}_{l}\}$ the serving cluster of user $k$, with
$K_{l}\triangleq\abs{\mathcal{D}_{l}}$ and $\abs{\mathcal{L}_{k}}$ its
cardinality. Full cell-free operation is the special case
$\mathcal{D}_{l}=\{1,\dots,K\}$ for all $l$, so $K_{l}=K$ and
$\abs{\mathcal{L}_{k}}=L$. A user outside $\mathcal{D}_{l}$ has
$\vect{w}_{kl}[q]=\vect{0}$, which leaves the signal model of
\eqref{eq:ap-transmit} untouched.

The frame timing, in contrast, is shared. Under \gls{tdd} the APs must switch
between directions together, since an AP still transmitting while a neighbour
receives creates cross-link interference between nodes that are far closer to
each other than to any user. The frame parameters
$\tau_{\mathrm{DL}},\tau_{\mathrm{UL}},\tau_{i,\mathrm{sig}}$ of
\Cref{sec:frame} are therefore network-wide constants, and only the load is
allowed to differ between APs. Repeating the construction of
\eqref{eq:load_avg} over the users AP $l$ actually serves gives the per-AP
average physical resource load
\begin{equation}
\label{eq:load_ap}
\bar x_{i,l} = \frac{N_{a,i}}{N_{i}}\,\frac{K_{i,l}}{K_{\max}},
\end{equation}
with $K_{i,l}=\abs{\mathcal{D}_{l}}$ counted over the users active in direction
$i$ and $K_{\max}$ the network-wide number of spatial streams. The
frame-averaging function \eqref{eq:frameavg} is reused with this per \gls{ap} load,
\begin{align}
\label{eq:frameavg_ap}
\mathcal{A}_{i,l}^{c}\!\left(P^{\mathrm a},P^{\mathrm s},P^{0}\right)
\triangleq{}&
\bar x_{i,l}\,\tau_i\,(1-\tau_{i,\mathrm{sig}})\,P^{\mathrm a}
+ \tau_i\,\tau_{i,\mathrm{sig}}\,P^{\mathrm s} \nonumber\\
&+ (1-\bar x_{i,l})\,\tau_i\,(1-\tau_{i,\mathrm{sig}})\,\delta^{\mu}_c\,P^{0}
\nonumber\\
&+ (1-\tau_i)\,\delta^{0}_c\,P^{0},
\end{align}
and the shorthand $\mathcal{A}_{i,l}^{c}(P)\triangleq
\mathcal{A}_{i,l}^{c}(P,P,P)$ is kept for the blocks whose consumption does not
depend on the operating mode. Under full cell-free operation all APs share the
same load, $\bar x_{i,l}=\bar x_{i}$, and \eqref{eq:frameavg_ap} collapses to
\eqref{eq:frameavg}.

\paragraph{Reconciling the two overhead accountings}
\todo{double check this part and work out some inconcistencies in the notation used above wrt prolog factors}
\Cref{sec:dmimo_sysmodel} charges the pilot overhead through the prelog
$(1-\tau_{p}/\tau_{c})$ of the coherence-block bookkeeping, while
\Cref{sec:rate} charges it through $\tau_{i}(1-\tau_{i,\mathrm{sig}})$. These are
two descriptions of the same time budget and must be identified rather than
multiplied. Under TDD reciprocity the pilots are sent in the uplink, so the
uplink signalling phase of \Cref{sec:frame} \emph{is} the pilot phase of the
coherence block, and writing $\tau_{d}^{\mathrm{DL}}$ for the number of samples
per block that carry downlink data, the identification is
\begin{equation}
\label{eq:overhead_map}
\tau_{\mathrm{UL}}\,\tau_{\mathrm{UL,sig}} = \frac{\tau_{p}}{\tau_{c}},
\qquad
\tau_{\mathrm{DL}}\,(1-\tau_{\mathrm{DL,sig}}) = \frac{\tau_{d}^{\mathrm{DL}}}{\tau_{c}},
\end{equation}
with $\tau_{d}^{\mathrm{DL}}\leq\tau_{c}-\tau_{p}$, the slack being the uplink
data and the downlink demodulation reference signals. The downlink-only prelog
$(1-\tau_{p}/\tau_{c})$ used in \Cref{sec:dmimo_sysmodel} is the special case
$\tau_{d}^{\mathrm{DL}}=\tau_{c}-\tau_{p}$, in which the whole block after the
pilots carries downlink data and neither uplink data nor downlink reference
signals are scheduled. The rate expression \eqref{eq:rate} then applies the
prelog once and must not apply it again. A second bookkeeping detail: the
subcarrier count is written $q_{i}$ in \eqref{eq:rate}, $Q_{\max}$ in
\eqref{eq:load_inst}, and $Q$ in \Cref{sec:dmimo_sysmodel}, while $q$ is also the
subcarrier index. We use $Q$ for the number of evaluated data subcarriers
throughout this section. \todo{align the symbols for the subcarrier count across
\Cref{sec:pwr_model} and \Cref{sec:dmimo_sysmodel}.}

\subsection{Per-AP power amplifiers}
\label{ssec:dmimo_pa}

Each AP is an independent transmitter with its own budget
\eqref{eq:power-constraint}, so the single quantity $P_{T,\mathrm{DL}}$ of
\Cref{sec:pa} becomes a vector of per-AP transmit powers. From the precoders of
\Cref{ssec:designs}, the power actually radiated by AP $l$ over the OFDM block is
\begin{equation}
\label{eq:ptl}
P_{T,l} \triangleq \sum_{q=1}^{Q}\sum_{k=1}^{K}
\expt{\norm{\vect{w}_{kl}[q]}^{2}} \leq \rho_{\max},
\end{equation}
which the two operation modes fill in differently. Local operation meets the
constraint with equality by construction, $P_{T,l}=\rho_{\max}$ for every AP,
because \eqref{eq:fractional-power} splits each AP's own budget among its users.
Centralized operation cannot rescale APs independently without breaking the
interference nulling, so the common factor of \eqref{eq:power-normalization}
gives $P_{T,l}=\gamma^{2}\tilde p_{l}$ with $\gamma^{2}=\rho_{\max}/\max_{l'}\tilde
p_{l'}$, and only the busiest AP reaches its budget. It is convenient to work
with the per-AP power utilisation
\begin{equation}
\label{eq:util}
u_{l} \triangleq \frac{P_{T,l}}{\rho_{\max}} \in [0,1].
\end{equation}

As was done in the centralized power model we assume that the per-AP power is shared equally across its antennas, the per-PA
output power is $P_{a,l}=P_{T,l}/M$ and the frame-averaged PA consumption of AP
$l$ follows \eqref{eq:pa_avg} with the per-AP load,
\begin{equation}
\label{eq:pa_avg_ap}
\overline{P_{\mathrm{PA},l}} = M\,
\mathcal{A}_{\mathrm{DL},l}^{\mathrm{PA}}\!\left(
P_{\mathrm{PA}}(P_{a,l}),\,
P_{\mathrm{PA}}(\zeta_{\mathrm{sig}}P_{\max}),\,
P_{\mathrm{PA}}(0)\right),
\end{equation}
with $P_{\max}$ the maximum PA output power and $\zeta_{\mathrm{sig}}$ the
fraction of it used during signalling, both as in \eqref{eq:pa_avg}.

Note that the equal-split assumption is not exact: the per-antenna powers
$\sum_{q}\sum_{k}\expt{\abs{[\vect{w}_{kl}[q]]_{m}}^{2}}$ differ across
$m=1,\dots,M$ for any precoder that is not isotropic. Because the efficiency
exponent $\alpha\in[0,1]$ of \eqref{eq:pa} makes $p\mapsto p^{\alpha}$ concave,
and hence $P_{\mathrm{PA}}(\cdot)$ concave, Jensen's inequality gives
$\frac{1}{M}\sum_{m}P_{\mathrm{PA}}(P_{a,lm}) \leq P_{\mathrm{PA}}(P_{a,l})$: the
equal-split assumption is a conservative upper bound on the true PA consumption,
and it is tight when the precoder loads the antennas of an AP uniformly.

\begin{remark}[PA sizing and the static floor]
\label{rem:pa_sizing}
The relation between $P_{\max}$ and the AP budget must be stated, because it
decides how the PA static term scales with the deployment. If each PA is
dimensioned for its share of the AP budget, $P_{\max}=\rho_{\max}/M$, then
substituting \eqref{eq:pa} into \eqref{eq:pa_avg_ap} makes the data-mode PA
consumption of AP $l$ independent of $M$,
\begin{equation}
\label{eq:pa_compact}
M\,P_{\mathrm{PA}}(P_{a,l}) =
\frac{\rho_{\max}}{\eta_{\mathrm{PA,max}}}
\Big[\xi + (1-\xi)\,u_{l}^{\alpha}\Big],
\end{equation}
where $\xi\in[0,1]$ is the static-to-dynamic consumption ratio and
$\eta_{\mathrm{PA,max}}$ the maximum PA efficiency of \eqref{eq:pa}. A fully
loaded AP therefore consumes exactly $\rho_{\max}/\eta_{\mathrm{PA,max}}$, and
the network static floor is $L\,\xi\rho_{\max}/\eta_{\mathrm{PA,max}}$, growing
with the number of APs rather than with the number of antennas. If instead
$P_{\max}$ is a fixed component rating independent of $M$, the floor is
$LM\,\xi P_{\max}/\eta_{\mathrm{PA,max}}$ and grows with the total antenna count.
The two conventions differ by a factor $M P_{\max}/\rho_{\max}$ and must not be
mixed. The three PA parameters $\xi$, $\alpha$, and $\eta_{\mathrm{PA,max}}$
were fitted for macro-cell amplifiers and should be refitted for the small
amplifiers of a distributed AP rather than reused.
\todo{source $\xi,\alpha,\eta_{\mathrm{PA,max}}$ for AP-class PAs.}
\end{remark}

\begin{remark}[Cost of the per-AP constraint]
\label{rem:gamma}
\todo{check this}
Under centralized operation with either nominal allocation of
\Cref{sssec:centralized}, the nominal powers satisfy $\sum_{l}\tilde
p_{l}=\sum_{k}\rho_{k}=L\rho_{\max}$, so the average utilisation is
\begin{equation}
\label{eq:ubar}
\bar u \triangleq \frac{1}{L}\sum_{l=1}^{L} u_{l}
= \frac{\gamma^{2}}{L\rho_{\max}}\sum_{l=1}^{L}\tilde p_{l}
= \gamma^{2}.
\end{equation}
The common rescaling factor of \eqref{eq:power-normalization} is therefore
exactly the fraction of the network power budget that is actually radiated.
Combining \eqref{eq:pa_compact} with \eqref{eq:ubar} and the concavity of
$u\mapsto u^{\alpha}$ bounds the network PA consumption in the data mode from
below by
\begin{equation}
\label{eq:pa_network_bound}
\sum_{l=1}^{L} M\,P_{\mathrm{PA}}(P_{a,l}) \geq
\frac{L\rho_{\max}}{\eta_{\mathrm{PA,max}}}
\Big[\xi + (1-\xi)\,\gamma^{2\alpha}\Big],
\end{equation}
with equality when all APs are equally loaded. The per-AP constraint thus costs
twice: the network radiates only a fraction $\gamma^{2}$ of $L\rho_{\max}$, yet
the static floor $L\xi\rho_{\max}/\eta_{\mathrm{PA,max}}$ is paid in full
whenever the PAs are not asleep. An imbalanced deployment, in which one AP is
much busier than the rest and drives $\gamma$ down, is penalised on both counts.
\end{remark}

\subsection{Per-AP analog front end}
\label{ssec:dmimo_ana}

A distributed AP carries few antennas, so hybrid beamforming has nothing to
amortise and the arrays are fully digital: $M_{\mathrm{RF}}=M$ and
$M_{\mathrm{PS}}=0$. Specialising \eqref{eq:ana_dl} and \eqref{eq:ana_ul}
accordingly, the working-mode analog consumption of one AP is
\begin{align}
\label{eq:ana_dl_ap}
P_{\mathrm{ana,DL},l} &= P_{\mathrm{LO}} + M\big(
2P_{\mathrm{DAC}} + 2P_{\mathrm{filt,RF}} + 2P_{\mathrm{mix}}\big),\\
\label{eq:ana_ul_ap}
P_{\mathrm{ana,UL},l} &= P_{\mathrm{LO}} + M\big(
2P_{\mathrm{ADC}} + 2P_{\mathrm{filt,RF}} + 2P_{\mathrm{mix}}
+ P_{\mathrm{LNA}}\big),
\end{align}
and the frame-averaged analog block of AP $l$ is
\begin{align}
\label{eq:ana_avg_ap}
\overline{P_{\mathrm{ana},l}} ={}&
\mathcal{A}_{\mathrm{DL},l}^{\mathrm{ana}}\!\left(P_{\mathrm{ana,DL},l}\right)
+ \mathcal{A}_{\mathrm{UL},l}^{\mathrm{ana}}\!\left(P_{\mathrm{ana,UL},l}\right)
\nonumber\\
&+ P_{\mathrm{ana,misc}} + P_{\mathrm{sync}} .
\end{align}
The local oscillator is the first term to change character. In the co-located
model a single LO is shared by $M_{\mathrm{ant}}$ chains and contributes a
constant; here every AP needs its own, so the network pays $L\,P_{\mathrm{LO}}$
and the term grows with the deployment even though it is small per node.

The term $P_{\mathrm{sync}}$ has no counterpart in \Cref{sec:ana} and is added
here. Coherent joint transmission across physically separate nodes requires the
APs to share a frequency and phase reference and requires the TDD chains to be
reciprocity-calibrated; the precoder of \eqref{eq:zf-precoder} nulls
interference only if the relative phases across APs are the ones the CPU assumed.
Whether the reference is distributed over the fronthaul, over a dedicated timing
network, or over the air, it costs power at every AP continuously, which is why
$P_{\mathrm{sync}}$ sits outside the direction-dependent averaging. It is a real
cost of distributing the array and one that a co-located design does not pay.
\todo{source a value for $P_{\mathrm{sync}}$; over-the-air calibration also
consumes frame time, which is not yet charged to $\tau_{i,\mathrm{sig}}$.}

\subsection{Functional split of Digital Signal Processing}
\label{ssec:split}

In the single-site model every digital operation runs at the base station. In a
distributed deployment the encoder, the precoder computation, and the precoder
application can each sit at the CPU or at the AP, and the choice sets both the
digital consumption and the fronthaul rate. We distinguish three splits, which
span the practical range and map onto the operation modes of
\Cref{ssec:designs}; \Cref{tab:split} summarises where each operation ends up.

\begin{table}[t]
\centering
\caption{Node at which each digital operation runs, per functional split. The
uplink combiner follows the same allocation as the downlink precoder. Splits S1
and S2 both realise the centralized precoding of \Cref{sssec:centralized} and
deliver the same rate; S3 is local operation (\Cref{sssec:local}) and delivers
the lower rate of local precoding. $^{\dagger}$Listed for completeness but not
priced, since \Cref{sec:dmimo_sysmodel} assumes perfect CSI \todo{add this}.}
\label{tab:split}
\ra{1.15}
\begin{tabular}{@{}llccc@{}}
\toprule
Operation & Count & S1 & S2 & S3 \\
\midrule
Encoding / decoding    & \eqref{eq:enc}  & CPU & CPU & CPU \\
Channel estimation$^{\dagger}$ & ---     & CPU & CPU & AP  \\
Precoder computation   & \eqref{eq:gops_cen}, \eqref{eq:gops_loc}
                                         & CPU & CPU & AP  \\
Precoder application   & $2K'N$          & CPU & AP  & AP  \\
OFDM (I)FFT            & \eqref{eq:ifft} & AP  & AP  & AP  \\
Predistortion          & \eqref{eq:dpd}  & AP  & AP  & AP  \\
Baseband filter        & \eqref{eq:bb}   & AP  & AP  & AP  \\
\bottomrule
\end{tabular}
\end{table}

Split~S1 forwards precoded samples. The CPU encodes, computes the precoder of
\eqref{eq:zf-precoder}, and applies it, sending AP $l$ the $M$ streams of
frequency-domain samples it must radiate; the AP performs only the OFDM
modulation, predistortion, baseband filtering, and conversion. 

Split~S2 forwards coefficients and data. The CPU encodes and computes $\mat{W}_{l}[q]$ but AP $l$
applies it locally, so the fronthaul carries the $K_{l}$ data streams plus the
precoding coefficients, refreshed once per coherence block. 


Split~S3 is local operation as defined in \Cref{sssec:local}: the CPU only encodes, and AP $l$
estimates its own channels and computes and applies the precoder of
\eqref{eq:local-rzf} from local information, so the fronthaul carries data alone.
Splits~S1 and~S2 both realise the centralized precoding of
\Cref{sssec:centralized} and therefore deliver the same rate; they differ only in
where the work is done and what crosses the fronthaul. Split~S3 delivers the
lower rate of local precoding.

\subsection{Digital signal processing at the APs and the CPU}
\label{ssec:dmimo_dig}

The per-chain digital costs of \Cref{sec:dig} carry over unchanged, evaluated
with $M_{\mathrm{RF}}=M$ at each AP. Every AP performs the OFDM modulation and
demodulation, the predistortion of its own PAs, and the baseband filtering
whatever the split, so
\begin{align}
\label{eq:dig_ap}
P_{\mathrm{dig,DL},l} &= P_{\mathrm{IFFT}} + P_{\mathrm{DPD}}
+ P_{\mathrm{filt,BB}} + P_{\mathrm{MIMO,DL},l},\\
P_{\mathrm{dig,UL},l} &= P_{\mathrm{IFFT}} + P_{\mathrm{filt,BB}}
+ P_{\mathrm{MIMO,UL},l},
\end{align}
with $\overline{P_{\mathrm{dig},l}}=
\mathcal{A}_{\mathrm{DL},l}^{\mathrm{dig}}(P_{\mathrm{dig,DL},l})+
\mathcal{A}_{\mathrm{UL},l}^{\mathrm{dig}}(P_{\mathrm{dig,UL},l})$. Only the
MIMO-processing terms $P_{\mathrm{MIMO},i,l}$ depend on the split, and they are
where the architecture actually shows.

As in \Cref{sec:dig}, a MIMO-processing term is a static part plus
$f_{s,\mathrm{I}}\Xi/\eta_{\mathrm{pre}}$, with $\Xi$ its complex operations per
sample. Two operations must be kept apart, because the split moves them
independently. Applying an already-computed precoder to $K'$ users over $N$
chains costs $2K'N$ operations per sample, the count already used for the
combiner in \Cref{sec:dig}. Computing it costs the Cholesky terms of
\eqref{eq:gops_pre}, amortised over the $\upsilon_{\mathrm{coh}}$ samples of a
coherence block. Under centralized operation the CPU inverts the $K\times K$ Gram
matrix of the collective channel $\mat{H}[q]\inC{LM\times K}$, so
\eqref{eq:gops_pre} applies with $M_{\mathrm{RF}}$ replaced by the total number
of distributed antennas $LM$,
\begin{equation}
\label{eq:gops_cen}
\Xi_{\mathrm{pre}}^{\mathrm{cen}} =
\frac{1}{\upsilon_{\mathrm{coh}}}\left(\frac{K^{3}}{3}
+ 3K^{2}LM + K\,LM\right),
\end{equation}
written here as a network total rather than per chain. Under local operation AP
$l$ instead inverts the $M\times M$ matrix of \eqref{eq:local-rzf}, which is the
same Cholesky-based accounting with the roles of the antenna and user dimensions
exchanged,
\begin{equation}
\label{eq:gops_loc}
\Xi_{\mathrm{pre},l}^{\mathrm{loc}} =
\frac{1}{\upsilon_{\mathrm{coh}}}\left(\frac{M^{3}}{3}
+ 3M^{2}K_{l} + M\,K_{l}\right).
\end{equation}
Setting $L=1$ and $M=M_{\mathrm{RF}}$ in \eqref{eq:gops_cen} recovers the
computation part of $M_{\mathrm{RF}}\Xi_{\mathrm{pre}}$ in \eqref{eq:gops_pre}
exactly.

The three splits then distribute these two counts over the two nodes as
\begin{align}
\label{eq:gops_split}
\Xi^{\mathrm{CPU}}_{i} &= \iota_{\mathrm{S1}}\,2K\,LM
+ \iota_{\mathrm{S1,S2}}\,\iota_{\mathrm{DL}}\,\Xi_{\mathrm{pre}}^{\mathrm{cen}},\\
\Xi_{i,l}^{\mathrm{AP}} &= \iota_{\mathrm{S2,S3}}\,2K_{l}M
+ \iota_{\mathrm{S3}}\,\iota_{\mathrm{DL}}\,\Xi_{\mathrm{pre},l}^{\mathrm{loc}},
\end{align}
where each $\iota_{\bullet}\in\{0,1\}$ equals one for the split or direction
named in its subscript and zero otherwise, and the corresponding powers follow
\eqref{eq:pre},
\begin{equation}
\label{eq:pmimo}
P_{\mathrm{MIMO},i,l} = \bar M_{\mathrm{RF}}P_{s}
+ \frac{f_{s,\mathrm{I}}}{\eta_{\mathrm{pre}}}\,\Xi_{i,l}^{\mathrm{AP}},
\qquad
\bar M_{\mathrm{RF}} = \Big\lceil \tfrac{M}{32} \Big\rceil ,
\end{equation}
and likewise for $P_{\mathrm{MIMO},i}^{\mathrm{CPU}}$ with
$\Xi^{\mathrm{CPU}}_{i}$. As in \Cref{sec:dig}, the computation is charged to the
downlink only, since TDD reciprocity lets the same matrix serve the uplink
combiner, while the application is charged in both directions. Every operation
then appears exactly once across \eqref{eq:gops_split}, so no work is counted
twice: under S2, in particular, the AP is billed only for applying the precoder
and the CPU only for computing it. Because linear combining is distributive over
the APs, $\hat s_{k}=\sum_{l}\vect{v}_{kl}^{H}\vect{y}_{l}$, an AP holding its
combining coefficients can form its own partial sum locally, so S2 attains the
centralized uplink rate while forwarding $K_{l}$ scalar streams instead of $M$
sample streams.

Summed over the network, \eqref{eq:gops_loc} grows linearly in $L$ at fixed $M$
while \eqref{eq:gops_cen} grows linearly in $LM$ and additionally carries the
$K^{3}$ term. For the dense, small-$M$ deployments of interest the centralized
cost is dominated by $3K^{2}LM$ and the local cost by $3M^{2}K_{l}$, so
centralized precoding is more expensive by roughly a factor $K/M$ per antenna.

Channel encoding and decoding stay at the CPU in all three splits, since that is
where the payload enters and leaves the network. The CPU consumption is
\begin{align}
\label{eq:cpu}
\overline{P_{\mathrm{CPU}}} ={}& P_{\mathrm{CPU},0}
+ \mathcal{A}_{\mathrm{DL}}^{\mathrm{dig}}\!\left(P_{\mathrm{enc}}
+ P_{\mathrm{MIMO,DL}}^{\mathrm{CPU}}\right) \nonumber\\
&+ \mathcal{A}_{\mathrm{UL}}^{\mathrm{dig}}\!\left(P_{\mathrm{dec}}
+ P_{\mathrm{MIMO,UL}}^{\mathrm{CPU}}\right),
\end{align}
with $P_{\mathrm{enc}}$ and $P_{\mathrm{dec}}$ from \eqref{eq:enc} driven by the
network sum rates $R_{\mathrm{DL}}$ and $R_{\mathrm{UL}}$, and
$P_{\mathrm{CPU},0}$ the always-on consumption of the central unit. The averaging
in \eqref{eq:cpu} uses the network-wide load $\bar x_{i}$, not a per-AP one.

\begin{remark}[Loss of static-power amortisation]
\label{rem:fpga}
The precoder, combiner, and baseband filter of \Cref{sec:dig} share a static
power $P_{s}$ per FPGA, one per $\bar M_{\mathrm{RF}}=\lceil
M_{\mathrm{RF}}/32\rceil$ group of chains. A co-located array of
$M_{\mathrm{ant}}$ antennas amortises that static power over 32 chains at a time;
a distributed network of $L$ APs with $M<32$ antennas each pays it $L$ times,
once per AP, for a total of $LM$ chains. At $L=40$ and $M=4$ the network runs
$40$ FPGAs where a co-located array of the same $160$ antennas would run $5$, an
eightfold increase in this static term. Chopping a large array into many small
ones leaves the per-chain digital costs untouched but destroys the sharing of
the fixed costs, and this is a general feature rather than an artefact of the
value $32$.
\end{remark}

\subsection{Fronthaul consumption}
\label{ssec:dmimo_fh}

Following the cell-free power model of \cite{ngo2018total}, the fronthaul link
serving AP $l$ consumes a traffic-independent part and a part proportional to the
rate it carries,
\begin{equation}
\label{eq:fh_basic}
P_{\mathrm{FH},l} = P_{\mathrm{FH},0} + \Pi_{\mathrm{FH}}\,R_{\mathrm{FH},l},
\end{equation}
with $P_{\mathrm{FH},0}$ in watts, the traffic-dependent coefficient
$\Pi_{\mathrm{FH}}$ in W per bit/s, and $R_{\mathrm{FH},l}$ the fronthaul rate in
bit/s. The static part covers the optical transceivers at both ends and depends
on the link length and the topology; the traffic part covers the
throughput-dependent consumption of the transport chain. In the parametrisation
of \cite{ngo2018total} these are $P_{\mathrm{FH},0}=\SI{0.825}{\watt}$ per link
and $\Pi_{\mathrm{FH}}=\SI{0.25}{\watt}$ per Gbit/s.
\todo{these two values are taken from a sub-6\,GHz cell-free study; check that
they are representative of the fronthaul technology assumed here before using
them for FR3 numbers.}

\paragraph{Fronthaul rate per split}
What $R_{\mathrm{FH},l}$ contains is set by the functional split of
\Cref{ssec:split}. Let $b_{\mathrm{FH}}$ be the number of bits per real sample
on the fronthaul, so a complex sample costs $2b_{\mathrm{FH}}$ bits. Under S1 the
link carries $M$ streams of frequency-domain samples at the constellation rate
$f_{s,\mathrm{I}}$, in both directions,
\begin{equation}
\label{eq:fh_s1}
R_{\mathrm{FH},l}^{\mathrm{S1}} = 2\,b_{\mathrm{FH}}\,M\,f_{s,\mathrm{I}},
\end{equation}
which becomes $2b_{\mathrm{FH}}Mf_{s,\mathrm{II}}$ if the CPU also performs the
OFDM modulation and time-domain samples are forwarded. Under S2 the link carries
the payload of the served users plus the precoding coefficients, an $M\times
K_{l}$ complex matrix per subcarrier refreshed once per coherence block of
$\upsilon_{\mathrm{coh}}$ samples,
\begin{equation}
\label{eq:fh_s2}
R_{\mathrm{FH},l}^{\mathrm{S2,DL}} = \sum_{k\in\mathcal{D}_{l}} R_{\mathrm{DL},k}
+ \frac{f_{s,\mathrm{I}}}{\upsilon_{\mathrm{coh}}}\,2\,b_{\mathrm{FH}}\,M\,K_{l}.
\end{equation}
Under S3 only the payload crosses in the downlink,
$R_{\mathrm{FH},l}^{\mathrm{S3,DL}}=\sum_{k\in\mathcal{D}_{l}}R_{\mathrm{DL},k}$.
In the uplink both S2 and S3 forward $K_{l}$ streams of per-AP partial sums
rather than $M$ streams of raw samples,
\begin{equation}
\label{eq:fh_ul}
R_{\mathrm{FH},l}^{\mathrm{S2,UL}} = R_{\mathrm{FH},l}^{\mathrm{S3,UL}}
= 2\,b_{\mathrm{FH}}\,K_{l}\,f_{s,\mathrm{I}},
\end{equation}
which for S2 costs nothing in rate, by the distributivity argument of
\Cref{ssec:dmimo_dig}. The power-control coefficients are refreshed only once per
large-scale fading realisation and are neglected, as in \cite{ngo2018total}; the
precoding coefficients of \eqref{eq:fh_s2} are not, because they change every
coherence block.

The signalling mode deserves separate treatment. Centralized operation requires
the CPU to see the raw received pilots in order to estimate the collective
channel, so during the uplink signalling phase the link must carry samples even
under S2, giving $R^{\mathrm{sig}}_{\mathrm{FH},l}=2b_{\mathrm{FH}}M
f_{s,\mathrm{I}}$ for S1 and S2 alike. Under S3 the pilots are consumed locally
and $R^{\mathrm{sig}}_{\mathrm{FH},l}\approx 0$. Centralization therefore imposes
a sample-level fronthaul during pilot transmission whatever the data-phase split.

\paragraph{Frame averaging}
The link is bidirectional and stays up across the whole frame, so its static term
is charged once per link with a single duty cycle rather than once per direction.
Let
\begin{align}
\label{eq:fh_duty}
\theta_{l} \triangleq \sum_{i\in\{\mathrm{DL},\mathrm{UL}\}}
\tau_{i}\Big[\tau_{i,\mathrm{sig}}
+ (1-\tau_{i,\mathrm{sig}})\,\bar x_{i,l}\Big]
\end{align}
be the fraction of the frame in which the link is in a working mode. With
$\delta^{\mu}_{\mathrm{FH}}$ the reduction factor of an idle transceiver, the
frame-averaged fronthaul consumption is
\begin{equation}
\label{eq:fh_avg}
\overline{P_{\mathrm{FH},l}} =
\Big[\theta_{l} + (1-\theta_{l})\,\delta^{\mu}_{\mathrm{FH}}\Big]
P_{\mathrm{FH},0}
+ \Pi_{\mathrm{FH}}\,\overline{R}_{\mathrm{FH},l},
\end{equation}
with the frame-averaged fronthaul rate obtained by reusing
\eqref{eq:frameavg_ap} with a zero base level, since an idle link carries no
traffic,
\begin{equation}
\label{eq:fh_rate_avg}
\overline{R}_{\mathrm{FH},l} = \sum_{i\in\{\mathrm{DL},\mathrm{UL}\}}
\mathcal{A}_{i,l}^{\mathrm{FH}}\!\left(
R^{i}_{\mathrm{FH},l},\,R^{i,\mathrm{sig}}_{\mathrm{FH},l},\,0\right).
\end{equation}
Optical transceivers wake slowly and hold most of their consumption while idle,
so $\delta^{\mu}_{\mathrm{FH}}$ is close to one in practice and the static
fronthaul term is nearly a constant per active link. This is what makes
$\mathcal{S}$, rather than the duty cycle, the effective lever on the static part.

\begin{remark}[Where the frame structure actually enters]
\label{rem:no_double}
The map from rate to power in \eqref{eq:fh_basic} is linear, and
$\mathcal{A}_{i,l}^{\mathrm{FH}}$ with a zero base is itself linear in its
arguments. For the data-sharing splits S2 and S3 the quantity to insert in
\eqref{eq:fh_rate_avg} is therefore the \emph{peak} rate during the data mode,
and since \eqref{eq:rate} already delivers $R_{\mathrm{DL},k}$ with the prelog
$\bar x_{i}\tau_{i}(1-\tau_{i,\mathrm{sig}})$ applied, the downlink data
contribution collapses back to $\sum_{k\in\mathcal{D}_{l}}R_{\mathrm{DL},k}$.
Charging the delivered rate directly is thus equivalent to frame-averaging the
peak rate, and doing both would count the frame structure twice. For the
sample-forwarding split S1 the situation is the opposite: $R^{\mathrm{S1}}_{\mathrm{FH},l}$
in \eqref{eq:fh_s1} is a hardware constant with no dependence on the delivered
rate at all, and the frame structure is then the only mechanism that reduces the
fronthaul load. The frame structure and the traffic dependence are therefore not
two independent knobs; each split makes one of them the active one.
\end{remark}

\begin{remark}[Fronthaul ceiling on energy efficiency]
\label{rem:fh_ceiling}
Under a data-sharing split the same payload is duplicated to every AP of a user's
serving cluster, so the network fronthaul load is
$\sum_{l}\sum_{k\in\mathcal{D}_{l}}R_{\mathrm{DL},k}=
\sum_{k}\abs{\mathcal{L}_{k}}R_{\mathrm{DL},k}$. Keeping only this term in
\eqref{eq:pnet} and using $\sum_{k}\abs{\mathcal{L}_{k}}R_{\mathrm{DL},k}\geq
(\min_{k}\abs{\mathcal{L}_{k}})\sum_{k}R_{\mathrm{DL},k}$ bounds the downlink
energy efficiency by
\begin{equation}
\label{eq:ee_ceiling}
\mathrm{EE}_{\mathrm{DL}} \leq
\frac{1}{\Pi_{\mathrm{FH}}\,\min_{k}\abs{\mathcal{L}_{k}}},
\end{equation}
independently of the bandwidth, the transmit power, and the precoder. Full
cell-free operation has $\abs{\mathcal{L}_{k}}=L$, so with $L=40$ and
$\Pi_{\mathrm{FH}}=\SI{0.25}{\watt}$ per Gbit/s the network cannot exceed
$\SI{100}{\mega\bit\per\joule}$ however well everything else is designed.
Shrinking the serving clusters raises the ceiling in direct proportion, which is
why AP selection is the dominant lever on cell-free energy efficiency and why
this term, absent from the co-located model, can decide the comparison on its own.
\end{remark}

\subsection{Delivered rate and network energy efficiency}
\label{ssec:dmimo_ee}

The per-user downlink rate follows from the SINR of \eqref{eq:downlink-sinr}
inserted into the ergodic rate expression \eqref{eq:rate},
\begin{align}
\label{eq:rate_user}
R_{\mathrm{DL},k} ={}& \bar x_{\mathrm{DL}}\,\tau_{\mathrm{DL}}\,
(1-\tau_{\mathrm{DL,sig}})\,\frac{\tilde B}{Q} \nonumber\\
&\times \expt{\sum_{q=1}^{Q}\log_{2}\!\left(1
+\mathrm{SINR}_{k}^{\mathrm{dl}}[q]\right)},
\end{align}
with $R_{\mathrm{DL}}=\sum_{k=1}^{K}R_{\mathrm{DL},k}$ the delivered sum rate.
The per-user decomposition, which the co-located model did not need, is required
here because the fronthaul term \eqref{eq:fh_s2} charges each AP only for the
users in its cluster. The network energy efficiency is then
\begin{equation}
\label{eq:ee_net}
\mathrm{EE}_{\mathrm{net}} =
\frac{R_{\mathrm{DL}} + R_{\mathrm{UL}}}{P_{\mathrm{net}}}
\quad[\mathrm{bit/J}],
\end{equation}
which reduces to \eqref{eq:ee} for a single AP without fronthaul. Because a
distributed and a co-located deployment cover the same area with different
hardware counts, a comparison at equal $L M$ and equal total radiated power is
the meaningful one, and the area energy efficiency
$\mathrm{EE}_{\mathrm{net}}/A$ over the served area $A$ is the fairer figure when
the deployments differ in extent.

A capacity constraint should accompany \eqref{eq:ee_net}. If the fronthaul links
have a finite capacity $C_{\mathrm{FH}}$, then $R_{\mathrm{FH},l}\leq
C_{\mathrm{FH}}$ must hold for every $l$, and under S1 this is a hard constraint
on $M$ and the sampling rate that is independent of the traffic. When it binds,
the rate of \eqref{eq:rate_user} is not achievable as written and the model
returns an optimistic energy efficiency.

\subsection{What the extension does not yet capture}
\label{ssec:dmimo_gaps}

Several effects are left out and are flagged here rather than silently absorbed
into the parameters. Fronthaul quantization is modelled only through its rate:
the $b_{\mathrm{FH}}$ bits per sample of \eqref{eq:fh_s1} add a distortion term
to the effective SINR of \eqref{eq:downlink-sinr} that is not present in
\Cref{sec:dmimo_sysmodel}, so S1 is charged for its fronthaul load but not
debited for the accuracy it loses. Channel estimation processing is likewise
unpriced: \Cref{sec:dmimo_sysmodel} assumes perfect \gls{csi}, so neither the
$LM\times K$ estimation at the CPU under centralized operation nor the $M\times
K$ local estimation under S3 appears in \eqref{eq:cpu} or \eqref{eq:dig_ap},
although both scale with the deployment. The synchronization and
reciprocity-calibration overhead is charged as a continuous power
$P_{\mathrm{sync}}$ but not as frame time, even though over-the-air calibration
consumes symbols that would otherwise carry data. Finally, the reduction factors
$\delta^{\mu}_{c}$ and $\delta^{0}_{c}$, the supply-and-cooling efficiencies, and
the PA parameters are all inherited from a macro base station and are the most
likely source of systematic error when they are applied unchanged to a
small outdoor AP. \todo{refit the AP-class hardware parameters, or report results
as a sensitivity range over them.}
